\chapter{Proposal Feedback and Iteration Plan}
\section{Proposal Feedback Summary}
The FYP committee approved the proposal (project F26-064) with minor modifications. In summary: all promised
components must be completed; the database and schema must be clarified; the need for sign-up/login must be explained
or a guest mode provided; the research gap, planned work and limitations must be stated clearly; external tools and
all claims must be specified and documented; the system should work across camera angles, lighting, board types and
piece materials (glass, wooden, plastic); occlusion must be handled using temporal data; and latency must be minimal.
(The project name in the committee record, ``Mentora'', is a placeholder; the product is Chesstor.)

\section{Proposal Action Register}
\begin{longtable}{p{0.05\textwidth}p{0.24\textwidth}p{0.31\textwidth}p{0.12\textwidth}p{0.18\textwidth}}
\toprule
No. & Feedback & Action taken / planned & Status & Evidence\\
\midrule
1 & Complete all promised components & Scope split into Iteration 1 (core) and Iteration 2 (app, tutor, upload, analytics); mapping in Ch.~9 & In progress & Sec.~\ref{sec:iter1}, \ref{sec:iter2}\\
2 & Clarify database and schema & Guest-first data design with optional cloud account; schema defined (Sec.~\ref{sec:data}) & Designed; \todo{team to confirm choice of DBMS} & Sec.~\ref{sec:data}\\
3 & Justify login or provide guest mode & Guest mode is the default; account only for cross-device history and analytics & Designed & Sec.~\ref{sec:data}\\
4 & State research gap, plan, limitations & Chapter~2 gap/limitations; Chapter~8 limitations & Done (draft) & Ch.~2, 8\\
5 & Specify external tools and document claims & External-component table; unverified claims marked \textbackslash todo & Partly done & Sec.~\ref{sec:deps}, Table~2.1\\
6 & Work across angles, lighting, boards, materials & Disturbance matrix in evaluation plan; real multi-board dataset planned & Planned & Sec.~\ref{sec:valplan}\\
7 & Handle occlusion with temporal data & Occlusion = missing evidence; temporal window; tested & Implemented (simulated) & Ch.~6--7\\
8 & Minimal latency & Latency measured per confirmation; frame-rate and window profiling planned & Partly measured (simulated) & Sec.~\ref{sec:results}\\
\bottomrule
\end{longtable}

\section{Iteration Strategy}
Development follows a vertical-slice, risk-first strategy: first prove that a physical move can be mapped to a legal
transition (the highest-risk assumption), then add robust recognition, then temporal reconstruction, then mobile
deployment, and finally engine analysis and UI polish. Work is integrated through GitHub with feature branches,
pull-request review, and automated CI (lint, tests, a frozen-seed regression gate); tagged versions are released
by the CD workflow.

\section{Iteration 1 Plan}\label{sec:iter1}
\textbf{Objective:} the core module with CI/CD. \textbf{Features:} calibration/rectification, per-square observation,
temporal buffer, legal-transition inference (UC-03 to UC-14), uncertainty and manual correction, Stockfish
wrapper, evaluation harness. \textbf{Deliverables:} tested package, CI/CD, simulated ablation results.
\textbf{Evidence:} repository, CI runs, test report (Ch.~6--7).

\section{Iteration 2 Plan}\label{sec:iter2}
\textbf{Objective:} real data and the live application. \textbf{Features:} real-board dataset and learned recogniser,
automatic board detection and tracking, Flutter live-tracking screen with overlay, PGN upload analysis, tutor
feedback and analytics, mobile profiling. \textbf{Evidence:} real-board metrics, device latency, demo.

\section{Work Breakdown and Responsibilities}
As planned in the proposal: Huzaifa --- CV pipeline and board reconstruction; Hiyam --- state tracking and move
inference; Madiha --- Stockfish and tutoring logic; everyone --- integration, testing and evaluation. Repository
folders map to these owners (\texttt{vision/}, \texttt{tracker.py}/\texttt{candidates.py}, \texttt{engine.py}).
\todo{Replace with actual per-member task log from the issue tracker.}

\section{Progress Against Plan}
Core module, front-end on synthetic images, harness and CI/CD are implemented. The Flutter application, real-board
data collection and tutor/analytics backend have not started.
\todo{Insert planned-vs-actual milestone table with dates.}

\section{Delays, Deviations and Recovery Plan}
\textbf{Deviation 1:} the core was implemented and tested in Python/laptop first, before the phone application,
to de-risk state reconstruction (as recommended by the implementation specification). \textbf{Deviation 2:} the POC
shown at the proposal defense was a UI mock-up; the mid-evaluation requires a live/functional interface, which is the
top priority of the next phase. \textbf{Recovery:} start the Flutter camera/overlay slice and real-data capture
immediately, in parallel with learned-recogniser work.
